%%%%%
%
\section{Correction de pics isol�s dans une s�rie de valeurs modulaires}

\itemstyle{Op\'erateur~:} \operfmt{correc\_pics\_angle\_auto\_unit}
\index{correc\_pics\_angle\_auto\_unit}

%%%%%OBJET%%%%
\itemstyle{Objet~:}
Les pics isol�s sont d�finis comme les �chantillons dont
les �carts au pr�c�dant et au suivant sont 
de signe oppos�s, et sup�rieurs (en valeur absolue)
� un certain seuil. 

Cet op�rateur d�tecte les pics dans une s�rie de
donn�es, et les remplace par une interpolation entre
les valeurs contigu�s.

Cet op�rateur est adapt� au cas des variables d�finies sur un domaine
p�riodique (arithm�thiqe modulaire), typiquement les angles : la 
transition entre 359 degr�s et 1 degr� n'est pas consid�r�e comme 
une discontinuit�.


%%%%%ARGUMENTS%%%%
\bigskip
%%%%%ENTREES%%%%
\begin{tabarg}
\multicolumn{3}{l}{\titretab {Entr\'ees}}\\
\hline
$X$ & Vecteur, \texttt{float} & S�rie de donn�es originelle \\
$S_0$ & Coeff., \texttt{float} & Seuil de d\'efinition d'un saut (m\^eme unit� que $X$, doit
\^etre positif et inf�rieur � la moiti� du domaine de p�riodicit�)\\
unit & Coeff., \texttt{string} & unit� de $X$ (degr�s ou radian)\\
\end{tabarg}

\medskip
%%%%%SORTIES%%%%
\begin{tabarg}
\multicolumn{3}{l}{\titretab {Sorties}}\\
\hline
$X_c$ & Vecteur, \texttt{float} & S�rie corrig�e (m\^eme unit\'e que $X$)\\
\end{tabarg}


%%%%%FORMULE%%%%
\itemstyle{Formule}

%% correc_pics_angle_auto_unit
L'\'echantillon $i$ est consid\'er\'e comme un saut si :


$|X[i] - X[i-1]| > S_0$ et $|X[i] - X[i+1]| > S_0$  et $ (X[i] - X[i - 1]) * (X[i] - X[i + 1]) > 0$


%%%%%AUTEUR%%%%
\itemstyle{Auteur~:}
CNRM/GMEI/TRAMM

%%%%%REFERENCE%%%%
\itemstyle{R\'ef\'erences~:}
